\subsubsection{P2：任务级调度中的整数规划、目标规划与启发式优化}

P2 将多 AMR 取送货调度写成带候选路径选择的任务级联合优化问题。其核心不是单独决定任务归属，而是在同一模型中同时处理“谁执行、按什么顺序执行、每段移动走哪条候选路径”，并把时间窗、电量安全、完工时间和运行成本纳入统一评价。因此，P2 集中体现了整数规划、目标规划、动态规划和启发式优化等课程知识点。

\paragraph{0-1 混合整数规划。}
P2 的任务指派变量 $x_{ki}$、首任务变量 $s_{ki}$、相邻任务变量 $u_{kij}$ 和路径选择变量 $z_{kabq}$ 都是 0-1 决策变量。$x_{ki}$ 描述任务归属，$s_{ki}$ 和 $u_{kij}$ 描述 AMR 内部任务链，$z_{kabq}$ 描述同一 OD 下候选路径的选择。通过这些变量，任务分配、任务排序和路径选择被统一表达为混合整数规划中的组合决策。

\paragraph{链式序列、时间窗与电量约束。}
P2 不只要求每个任务被分配，还要求每台 AMR 的任务形成从初始位置出发的线性执行链。时间衔接约束保证后继任务必须在前序任务完成并空驶到取货点之后才能开始，释放时间和最晚完成时间用于刻画时间窗与延期。电量约束沿任务链递推，将空驶、载货和服务耗电计入剩余电量，保证任务级计划在路径和能源层面可执行。

\paragraph{分层目标规划。}
P2 采用目标规划中的优先级因子思想，而不是把所有指标简单相加。评价顺序先处理缺失路径和电量违反，再处理高优先级延期与延期任务数，随后比较时间效率 $F_2$，最后比较运行成本 $F_3$。这种字典序结构体现了 $P_1\gg P_2\gg P_3$ 的管理含义，避免用低层运行成本收益抵消高层服务质量或可行性损失。

\paragraph{分支定界、割平面与动态规划。}
m1 采用两阶段分解：阶段一是 0-1 指派模型，可对应分支定界思想；阶段二对每台 AMR 的车内任务排序使用 Held-Karp 型动态规划，状态由“已完成任务集合 + 当前末任务”构成。m2 则把指派、排序、路径选择、时间衔接和目标规划统一写入联合 MILP，通过分支定界和割平面思想搜索满足分层目标的解。两者共同体现了从分解建模到联合精确建模的运筹学求解路线。

\paragraph{启发式构造与局部搜索。}
m3 面向规模扩展，先用 regret-2 后悔值插入构造初始任务链，再通过 relocate、swap 和 reroute 三类邻域进行局部搜索。该方法不能证明全局最优，但在实验中能以远低于 m2 的时间得到相同或接近的解质量，体现了启发式优化在批量实验、较大规模算例和后续动态重排中的工程价值。